\clearpage
\section{Conclusion}
\label{sec:conclusion}

This thesis set out to address a fundamental tension in confined-space drone operations: when collisions with the environment are inevitable, how should the platform be designed to survive them, and can the onboard sensor signatures of those collisions be exploited for spatial awareness? The work followed two parallel threads --- mechanical deflection and inertial inference --- and the results from each thread reinforce a common theme: the rotating cage mechanism fundamentally redistributes collision energy, with measurable and operationally meaningful consequences for both survivability and sensor interpretation.

\subsection{Summary of Contributions}

\textbf{Platform Design.} A custom 4-inch quadcopter was designed, built, and flight-validated. The platform's dimensions were governed by a three-way trade-off: a lower bound set by the payload requirements of the Raspberry~Pi~5 companion computer and PETG protective cage, and an upper bound set by indoor safety limits on propeller size and kinetic energy. A 6S power system paired with 2004-stator motors provides a thrust-to-weight ratio of approximately 2.57:1 (target was 2.87:1 at 1,200\,g; final assembly mass $\sim$1,340\,g), placing hover at roughly 40\% throttle and reserving headroom for post-collision recovery. The autonomous flight stack integrates PX4 firmware, ROS~2 middleware, and an OptiTrack MoCap ground-truth system through the EKF2 estimator, with the \texttt{flight\_director.py} mission executor driving a reproducible collision-loop state machine. The PETG spherical cage (35.8\,cm diameter) was tested in two configurations --- rigidly mounted (Fixed Cage) and bearing-decoupled (Rotating Cage) --- isolating the bearing mechanism as the single experimental independent variable.

\textbf{Mechanical Performance.} Across 127 verified structural impacts (67 Fixed, 60 Rotating), the rotating cage reduced maximum deceleration by 37.2\% relative to the fixed configuration ($-1.64 \pm 0.64$\,m/s$^2$ vs.\ $-1.03 \pm 0.50$\,m/s$^2$, $p < 0.001$, $d = 1.05$). IMU vibration analysis revealed qualitatively distinct post-impact behaviors: the fixed cage sustained high-frequency pitch-and-roll oscillations consistent with stick-slip contact dynamics, while the rotating cage's bearing decoupling allowed the internal IMU to return to a nominal noise floor substantially faster. The aggregate IMU envelope (mean~$\pm\sigma$ across all 127 flights) confirmed a narrower acceleration-deviation spread for the rotating configuration.

The deflection benefit carries a modest spatial cost. The rotating cage exhibited a slightly wider maximum trajectory deviation ($131.0 \pm 32.1$\,mm vs.\ $118.4 \pm 30.4$\,mm, $p = 0.022$, $d = 0.41$), consistent with the bearing mechanism transferring angular momentum at the moment of contact. However, the total recovery area --- the Spatial Integral of Absolute Error --- remained statistically indistinguishable between configurations ($p = 0.998$), indicating that the flight controller's PID correction capacity absorbs this offset without expanding the overall spatial footprint of the recovery maneuver.

\textbf{Impact-Angle Inference.} A Random Forest regressor trained on 14 IMU-derived features recovered impact angle with $R^2 \approx 0.72$ and mean absolute error of $4.5^\circ$--$4.9^\circ$ for both cage configurations. This level of accuracy supports rudimentary tactile awareness --- the drone can distinguish a shallow glancing impact ($\sim$30$^\circ$) from a deeper collision ($\sim$60$^\circ$) using only the vibration signature recorded by its onboard IMU.

Three subsidiary findings carry practical significance. First, the rotating cage's angle signature is invisible to any linear model (Huber baseline $R^2 = 0.004$) yet fully recovered by the non-linear Random Forest ($R^2 = 0.719$), demonstrating that the bearing mechanism disperses angle information across a distributed multivariate fingerprint rather than destroying it. Second, cross-condition transfer fails decisively in both directions ($R^2 < 0.26$), meaning a model trained on one cage type cannot be deployed on the other without a mechanical-normalisation transfer function. Third, persistent learning-curve gaps of approximately 0.10 $R^2$ between training and validation performance indicate that the current dataset of 127 flights is insufficient to fully saturate the 14-dimensional feature space, and that additional data collection or stronger regularisation would close this gap.

\subsection{Future Work}

Two directions emerge naturally from the findings, and a third follows from exploratory work conducted alongside the main experiments.

\textbf{Expanding the collision envelope.} The experimental matrix should be extended to additional obstacle geometries (flat walls, convex corners, pipes of varying diameter), surface materials (metal, rubber, glass), and approach angles spanning the full $0^\circ$--$90^\circ$ range. Multi-impact sequences --- where the drone ricochets between surfaces before recovering --- would test whether the vibration-decoupling benefit of the rotating cage compounds across successive contacts.

\textbf{From Lab Angle Inference to Tactile SLAM.} The $\sim$5$^\circ$ angle-prediction accuracy demonstrated in this work was obtained under controlled laboratory conditions with MoCap ground truth. The natural next step is to progressively relax this dependency and move toward fully onboard inference. A practical intermediate milestone is to replace the external OptiTrack system with onboard visual-inertial odometry (VIO), comparing its velocity estimates against the EKF2 states used here, and confirming that the Random Forest angle predictions remain stable when fed with onboard-derived features rather than MoCap-derived ground truth. Beyond this, the platform should be exposed to genuinely unknown environments --- unstructured industrial mock-ups with feature-poor surfaces, multi-path radio reflections, and ferromagnetic interference --- where no external tracking is available and the drone must rely entirely on its onboard sensor suite. In such settings, impact-angle inference would serve as the primary input to an autonomous navigation policy: the drone would classify each collision by angle and severity, update an internal belief map of obstacle locations, and adapt its flight path in real time without ever requiring visual line-of-sight to the obstacle. If successive impacts against the same surface can be triangulated from their IMU signatures, a coarse geometric map of the environment can be built purely from contact signatures, realizing \emph{tactile simultaneous localization and mapping}. Combined with the mechanical resilience of the rotating cage, this closes the loop from the present work's \emph{post-hoc} angle estimation to a fully \emph{online} tactile navigation capability --- from a drone that is blind but survivable to one that is blind but tactile, navigating by touch in spaces where no other sensor modality is viable.

\textbf{Shell-based wall odometry.} Lerche~\cite{lercheSpinFlyUAVSystem} explored whether the rotating protective shell itself could serve as a contact-odometry sensor: a magnetic encoder coupled to the shell measured rotational displacement as the drone rolled along a wall, providing a direct estimate of traversal distance without relying on visual features or external tracking. A preliminary replication of this concept was attempted with the platform described in this work, but two mechanical limitations prevented reliable wall-distance tracking. First, the drone did not actively maintain wall contact---the rigid cage frame repeatedly deflected the vehicle away from the surface, producing intermittent contact rather than sustained rolling. Adding a rubber or high-friction band around the cage exterior could mitigate this by increasing surface grip and damping the bounce-back. Second, a quadcopter tilts forward during forward flight, meaning the outer ring of the cage is not horizontally aligned with a vertical wall surface; only a small arc of the cage rim makes contact, reducing the effective rolling circumference and introducing geometric error into the odometry reading. A gimbal-mounted cage that remains parallel to the wall surface regardless of the drone's pitch attitude would resolve this alignment problem, analogous to the rotating sensor head on the Elios series of inspection drones~\cite{briod2014collision}. \\[0.5em]
\begin{figure}[htbp]
    \centering
    \includegraphics[width=\textwidth]{Figures/wall_pass_drone.png}
    \caption{Preliminary wall-contact odometry test. The drone attempts to roll along a vertical surface using the rotating cage as a tactile sensor; intermittent contact and pitch-induced cage misalignment prevented reliable distance tracking in this early trial.}
    \label{fig:wall_pass}
\end{figure}
